\section{Workload Descriptions}
\label{sec:workloads}

UsingMatrixClass supplies four workload pairs.
Each pair is one shared header and two thin \texttt{main} wrappers.
The regular wrapper does nothing at start-up, so every \texttt{new} goes to
the C library.
The overloaded wrapper links the allocator's global \texttt{new} target.
It initializes the memory pool before the workload starts and shuts that pool
down after the workload ends.
Nothing else differs between the two programs of a pair.
The descriptions below are taken from the workload headers in the
UsingMatrixClass source tree.

\subsection{Properties shared by every workload}

\paragraph{Deterministic input.}
No workload reads the clock or the environment.
The matrix workloads fill matrices from a fixed deterministic generator, and
the linked list draws every decision from a generator with a fixed seed.
The regular and overloaded programs therefore execute the same sequence of
allocations, frees, and arithmetic.
Only the allocator differs.

\paragraph{Checksum.}
Every workload folds its results into a checksum and prints it on completion.
Matching checksums between the regular and overloaded runs prove that both did
the same work and that the overloaded allocator returned usable memory.
All 240 v6 and v7 runs and every M3~Max matrix row report matching checksums
(\Cref{sec:statistical-results}).

\paragraph{Pool sizing.}
The overloaded program has to tell the allocator how much memory to reserve
before the workload starts.
This is because the pool is one fixed \texttt{mmap} region that never grows.
Each workload computes that number from its own arguments.
The matrix workloads use the matrix size and count, the linked list uses its
maximum length, and the stress workload uses a fixed 512~MiB.
The regular program ignores the number.
One consequence matters when reading the results.
The overloaded program's peak memory reflects the pages it touched
inside that reservation, not the size of the reservation itself.

\paragraph{Single thread.}
Each workload program runs on one thread, and none of them creates additional
threads.
The allocator versions from v4 onward are built to be used by many threads at
once.
They add per-thread caches and lock-free ownership checks
(\Cref{sec:allocator-versions}).
However, the allocator does not create threads either.
In these experiments, therefore, the pool mutex is never contended and every
thread cache has exactly one owner.
What v4--v7 save here is the cost of the central bin search on a cache hit,
not time spent waiting for a lock.
The allocator's behavior under real concurrency is measured separately with
the multi-threaded matrix programs in \Cref{sec:pthreads-experiments}.

\subsection{Matrix array}

The matrix array workload lives in
\texttt{Matrix\allowbreak Array\allowbreak Workload.hpp}.
It creates 256 matrices of the same size $n$ and keeps them in one array for
the whole run.
The only command-line argument is the matrix size.
A matrix in MatrixClassDemo allocates each of its rows separately.
As a result, building one matrix costs $n$ allocations of about $8n$ bytes
each.

The program then runs 640 replacement passes.
During a pass, every matrix is replaced by a new matrix computed from the one
before it in the array.
Most of the time that computation is an addition.
For the 8 matrices whose index is a multiple of 32, it is a multiplication
instead.
Each computation produces a temporary result matrix, which is then copied into
the array slot and destroyed.
In other words, every pass allocates and frees 256 temporary matrices, and each
of those is $n$ row allocations.

The important thing about this workload is the balance between arithmetic and
allocation.
The allocations all have the same shape, so the allocator sees the same request
size over and over.
Meanwhile, the 8 multiplications in every pass each cost on the order of
$n^{3}$ floating-point operations, and they take up most of the run time.
It tests whether an overloaded allocator slows down a program that barely
needs it.
It is also the workload most sensitive to how rows are laid out in memory, which
is the reason the v8 redesign targeted it.

\subsection{Uniform matrix nodes}

The uniform nodes workload lives in
\texttt{Uniform\allowbreak Matrix\allowbreak Nodes.hpp}.
It also uses 256 matrices of one size $n$.
The difference is that each matrix is owned by its own separately allocated
node instead of sitting in a shared array.
The arguments are the matrix size and the number of replacement passes.

Each pass works the same way as the matrix array.
Every node is replaced by a new one computed from the node before it, using
addition most of the time and multiplication for every 32nd node.
However, replacing a node now means destroying the old node object and
allocating a brand new one to hold the result.
Because of that, the node's address changes on every pass, and each object
lives for only one pass rather than the whole run.

Since every matrix has the same size, the allocator sees only a few request
sizes.
These are the node object itself and the $n$ rows of its matrix.
Those sizes are requested and released constantly.
This is exactly the pattern that the v4--v7 thread cache and the v6--v7 slab
refills were designed for.
A block freed at one size is very likely to be requested again at that same
size a moment later.
Comparing this workload with the matrix array therefore isolates one thing.
The arithmetic is identical, so any difference between the two comes from the
added cost of destroying and recreating the node objects.

\subsection{Nested matrix stress}

The stress workload lives in
\texttt{Nested\allowbreak Matrix\allowbreak Stress.hpp}.
Instead of one flat array, it builds two trees of matrices.
There are six three-level trees (\texttt{Matrix***}) and four four-level trees
(\texttt{Matrix****}).
Each leaf of a tree holds between 6 and 18 matrices.
Unlike the other matrix workloads, the matrices here are not all the same
size.
Their dimensions are picked from twelve values ranging from 16 to 256.
That means row sizes range from 128 to 2048 bytes and whole matrices from
2~KiB to 512~KiB.
The only argument is a round count.

The leaf matrices are grouped in threes.
In each round, the program visits every group and stores the result of an
operation on the first two matrices into the third.
The operation is an addition most of the time and a multiplication every third
time.
Each operation allocates a fresh result matrix.
After the arithmetic, the round picks one outer branch of each tree at random,
destroys the whole branch, and rebuilds it with new random shapes.

This is the fragmentation workload.
Destroying a branch frees a large number of blocks of mixed sizes all at once.
The rebuild then immediately asks for a different mix of sizes.
That pattern puts pressure on the free-block data structure, which is a single
list in v1--v2 and size bins from v3 onward.
It also exercises the rule for splitting a block and returning the leftover.
Most of all, it separates the two coalescing policies.
Versions v1--v4 merge neighboring free blocks right away, while v5--v8 delay
merging until an allocation fails.
Finally, peak memory on this workload shows internal fragmentation.
The allocator rounds requests up to a 64-byte quantum, and most of the twelve
row sizes do not land on a quantum boundary, so some space is wasted in every
block.

\subsection{Linked list}

The linked list workload lives in
\texttt{Linked\allowbreak List\allowbreak Workload.hpp}.
It maintains a doubly linked list of \texttt{Node} objects.
A node is 72 bytes and holds two links, an index, a count, a pointer, six
integers, three characters, and four booleans.
Each node also owns a separately allocated array of 5 to 10 integers.
The only argument is the maximum list length, and the program performs 200
operations for every node of that maximum.

Every operation is chosen at random from a fixed mix.
Twenty percent of operations append a node.
Another twenty percent insert a node before a random existing node.
Twenty percent erase a random node.
Twenty-five percent update the fields of a random node without allocating.
The remaining fifteen percent replace a random node's integer array with a new
one of a different random length.
Once the list is full, append and insert turn into erase instead.
As a result, the list stays close to its maximum length for most of the run.

This workload is different from the other three because it does no matrix
arithmetic at all.
Each operation is a short pointer walk plus, in four of the five cases, one or
two allocations or frees of a few dozen bytes.
That makes it the pure allocator benchmark.
Run time is dominated by \texttt{new}, \texttt{delete}, and the cache misses
from touching millions of small objects scattered through memory.
It is also the workload where the allocator's own bookkeeping is proportionally
largest.
The header, back-pointer, and canaries described in
\Cref{sec:allocator-versions} add up to 64 bytes, which is almost as much as
the 72-byte node they are attached to.

\subsection{Run scales and parameter selection}

The assignment defines three run scales by how long they last.
Short runs take 3 to 5 minutes, medium runs take 2 to 3 hours, and long runs
take more than 6 hours.
Each scale requires five regular runs and five overloaded runs.
The parameters for each scale were chosen through calibration runs on each
machine, and those calibration runs do not count as measurements.
Because machines differ in speed, the same scale means different parameters on
different hosts.
\Cref{tab:workload-parameters} records the values that were run in every
campaign.

There were two different ways to make a run last longer, and they answer
different questions.

The first way is to repeat the same work.
The Group~2 campaigns on the Apple M3~Max and the Ryzen 7900X kept the program
arguments fixed and launched the executable several times in a row.
All of those launches were timed together as one run using
\texttt{repeat\allowbreak\_workload.sh}.
For example, a medium matrix run on the Ryzen host is \texttt{regular\_new 360}
launched 42 times in sequence, and a long run launches it 117 times.
In this approach the working set, the request sizes, and the pool size are the
same at every scale.
Only the number of repetitions changes.
As a result, these campaigns show almost no change in peak memory across
scales, and the gap between regular and overloaded \texttt{new} stays flat from
short to long.
One side effect is that start-up costs, such as the pool \texttt{mmap} and the
first-touch page faults, are paid once per repetition rather than once per
run.

The second way is to make the problem bigger.
The v2 campaign on the Core i7-1165G7 and the Group~1 campaigns launched each
executable once per measurement with a larger argument at each scale.
For example, the v2 matrix runs used sizes of 275, 800, and 1225, and the v2
linked list runs used maximum lengths of 125\,000, 580\,000, and 875\,000.
Here the working set grows along with the run time.
A larger matrix changes the row size and therefore the allocator's size class.
A longer list means more live blocks for a free-list search to walk through.
Peak memory grows as well.
This approach shows how an allocator responds to problem size.
However, a change from short to long cannot be blamed on duration alone.

These two approaches should not be compared with each other.
The flat scale behavior seen in the v6 and v7 campaigns should not be expected
in the v2 or Group~1 results.
Likewise, the very large v2 regressions on the stress and linked list
workloads at the long scale reflect a much bigger problem, not a longer clock.
\Cref{sec:statistical-results} keeps the two approaches in separate tables for
this reason.

\begin{table*}[!t]
\centering
\caption{Workload parameters that were run, by campaign. ``Reps'' is the
number of back-to-back executions timed as one run under the repeat strategy.
Values are from the campaign CSVs and \texttt{environment.txt} files; Group~1
values are from the shared data-collection sheet
(\Cref{sec:environments}). A dash means the scale was not run or not
recorded.}
\label{tab:workload-parameters}
\footnotesize
\setlength{\tabcolsep}{4pt}
\begin{tabularx}{\textwidth}{@{}l l l L L L l@{}}
\toprule
Campaign (host) & Versions & Workload & Short & Medium & Long & Strategy \\
\midrule
Group 2, Apple M3 Max & v1--v7 & Matrix array & 360, 1 rep & 360, 37 reps & 360, 97 reps & repeat (3 runs) \\
Group 2, Apple M3 Max & v6 & Matrix array & 360, 1 rep & 360, 37 reps & 360, 97 reps & repeat \\
 & & Uniform nodes & 350 700, 1 rep & 350 700, 37 reps & 350 700, 97 reps & repeat \\
 & & Nested stress & 800, 1 rep & 800, 34 reps & 800, 93 reps & repeat \\
 & & Linked list & 6\,500\,000, 1 rep & 6\,500\,000, 32 reps & 6\,500\,000, 84 reps & repeat \\
Group 2, Ryzen 9 7900X (WSL) & v7 & Matrix array & 360, 1 rep & 360, 42 reps & 360, 117 reps & repeat \\
 & & Uniform nodes & 350 700, 1 rep & 350 700, 42 reps & 350 700, 115 reps & repeat \\
 & & Nested stress & 800, 1 rep & 800, 38 reps & 800, 102 reps & repeat \\
 & & Linked list & 3\,512\,841, 1 rep & 3\,512\,841, 43 reps & 3\,512\,841, 120 reps & repeat \\
Group 2, Apple M3 Max & v8 & Matrix array & 360 & 360 & 360 & repeat \\
Group 2, Core i7-1165G7 (Ubuntu) & v2 & Matrix array & 275 & 800 & 1225 & grow \\
 & & Uniform nodes & 285 600 & 900 600 & 1200 600 & grow \\
 & & Nested stress & 100 & 440 & 900 & grow \\
 & & Linked list & 125\,000 & 580\,000 & 875\,000 & grow \\
Group 2, ThinkPad X13 (WSL) & v3 & Uniform nodes & 225 600 & 800 600 & -- & grow \\
 & & Nested stress & 200 & 1000 & -- & grow \\
Group 1, Ryzen 7 7800X3D (WSL2) & v4, v5 & Uniform nodes & 100 35\,000 & 100 1\,335\,000 & 100 4\,100\,000 & grow \\
Group 1, Apple M5 Pro & v4, v5 & Uniform nodes & 100 35\,000 & 100 1\,335\,000 & 100 4\,100\,000 & grow \\
Group 1, Core i5-6400 (Arch) & v5 & Matrix array & 400 & 900 & -- & grow \\
 & & Uniform nodes & 400 500 & -- & -- & grow \\
Group 1, Core i7-11800H & v5 & Matrix array & -- & -- & -- & not recorded \\
\bottomrule
\end{tabularx}
\end{table*}

\subsection{Expected sensitivity of each workload}

Reading the workload headers alongside the allocator versions in
\Cref{sec:allocator-versions} suggests where each pair should show a
difference and where it should not.
These expectations are written down here so that the results in
\Cref{sec:statistical-results,sec:charts-comparisons} can be judged against
them.

\begin{itemize}
  \item \textbf{Matrix array} should show the smallest effect of any version
    in either direction.
    The 8 matrix multiplications in each pass cost on the order of $n^{3}$
    operations each, and they swamp the 256 temporary-matrix allocations.
    Any difference that does appear is more likely to come from how rows are
    placed in memory than from allocation speed.
  \item \textbf{Uniform nodes} should reward versions that quickly reuse a
    freed block of the same size.
    That favors the v4--v7 thread cache and the v6--v7 slabs.
    It should penalize versions with an expensive free path, such as the
    v1--v2 address-ordered insertion with immediate merging.
  \item \textbf{Nested stress} should separate the two coalescing policies.
    Immediate merging in v1--v4 keeps large free spans available for the next
    rebuild.
    Deferred merging in v5--v8 makes each free cheap but risks a failed search
    followed by a full-region merge.
    This workload should also show the largest internal-fragmentation penalty
    in peak memory.
  \item \textbf{Linked list} should be the hardest workload for every custom
    version.
    The 64-byte header plus alignment and quantum rounding roughly triples the
    footprint of a 72-byte node.
    Meanwhile, the C library's small size classes and per-thread cache are
    already tuned for exactly this case.
\end{itemize}

Whether the measurements bear these expectations out is the subject of the
next two sections.
